\chapter{Marco Teórico}
\label{ch:marco}

\section{Descripción}

Toda tesis hace referencia a trabajos previos en el área y trabajos afines que
están directamente relacionados con lo planteado en la tesis.

Además, en el marco teórico debe aparecer la información absolutamente
necesaria para comprender la solución, y por eso es recomendable escribir
primero la solución (el siguiente capítulo), para ir anotando qué debe ser
explicado en el marco teórico.

\subsection{Redacción}

La \nt{redacción} en todo el documento debe seguir un estilo científico
objetivo. Esto implica que se redacta de modo impersonal, sin utilizar primeras
personas del singular o del plural, y se evita el uso de cualquier tipo de
calificativo, sustituyéndolos siempre por datos concretos, vinculados a
referencias bibliográficas o datos experimentales. Los comparativos también
deben concretarse a hechos y datos, y nunca dejarse ``en el aire''. Por la
naturaleza de la tesis, el tiempo verbal es usualmente presente, no perdiendo
nunca de vista que se está explicando ``cómo hacer algo'', en vez de ``qué se
hizo''.

Las \nt{frases} deben ser cortas, y debe evitarse que el lector tenga que saltar
constantemente entre partes de la tesis, lo que implica una exposición lineal
clara, donde lo que se necesita ya ha sido explicado antes. Deben evitarse
redundancias y por tanto cada concepto se exponen en un único lugar.

Todo aspecto circunstancial es irrelevante para la tesis, es decir, si se ha
desarrollado en el laboratorio $X$, o en el curso $Y$, con el profesor $Z$, o
en la empresa $W$, el nombre de funciones o clases en su código, etc., es
información irrelevante para reproducir el experimento, y por lo tanto sobra.
Numeración del documento

La primera página de la tesis es la correspondiente a la introducción, así que
ésta debe ser la página 1. Desde la introducción, hasta antes de la
bibliografía, las unidades son ``Capítulos''. La bibliografía y anexos no se
consideran capítulos, así que ya no continúan con la misma numeración de los
capítulos (la paginación sí continua). Los índices, notación, glosario, etc. se
numeran con números romanos en minúscula (i,ii,iii,iv,v,vi...) y antes del
índice (portada, resúmenes, agradecimientos, hoja de evaluadores, etc.) las
páginas no llevan numeración. 

Esta plantilla LaTeX ya se ocupa de todo lo anterior.

\subsection{Ecuaciones}

Para citar \nt{ecuaciones} se utilizan paréntesis redondos, y no es necesario emplear la palabra ``ecuación''. Por ejemplo ``Introduciendo en (4.2) los resultados de (3.3) y (3.7) se obtiene ...''. La ecuación es parte del flujo de texto y no un objeto flotante, así que no pueden emplearse como figuras. Cuando se requiere la ecuación, allí se inserta \eqref{eq:QP}.

\begin{equation}\label{eq:QP}
    \begin{split}
     \min_{\boldsymbol{x}_{k+1}, \boldsymbol{u}_{k}} \quad \boldsymbol{J}&= \left (\boldsymbol{x}_{k+1} - \boldsymbol{x}^{*}\right )^{T}{\boldsymbol{Q}}\left (\boldsymbol{x}_{k+1} - \boldsymbol{x}^{*}\right )\\&+\left ( \boldsymbol{u}_{k}- \boldsymbol{u}^{*} \right )^{T}\boldsymbol{R}\left ( \boldsymbol{u}_{k}- \boldsymbol{u}^{*} \right ) \\  \mathrm{s.a.} \quad \boldsymbol{x}_{k+1}&={\boldsymbol{A}}\boldsymbol{x}_{k} + {\boldsymbol{B}}_{k} \boldsymbol{u}_{k} + {\boldsymbol{d}}_{k} \\
     \boldsymbol{G}\boldsymbol{u}_{k} &\geq \boldsymbol{W}
    \end{split}
\end{equation}

También puede utilizar el editor de ecuaciones en línea.

\subsection{Figuras}

Para el almacenamiento de imágenes existen dos tipos de formato: las imágenes
raster y las imágenes vectoriales.\index{imagen!raster}

Las imágenes raster son representadas por una rejilla de píxeles, en donde cada
píxel tiene un valor que representa al nivel de gris o el color. La
discretización espacial es ineludible, y la única forma de obtener buena
calidad es empleando tamaños grandes de la imagen que conduzcan a resoluciones
de al menos 300 puntos por pulgada en la impresión, lo que conlleva a archivos
de documentos de varios megabytes. Dentro de los formatos para almacenar
imágenes raster existen algunos con pérdida (como el JPEG) que producen en
imágenes sintéticas, como diagramas, estructuras ruidosas que dan una
apariencia de baja calidad a las figuras. Otros formatos (como PNG, BMP, TIFF o
GIF) no tiene pérdidas de información, pero los algoritmos de compresión no
pueden reducir el tamaño de las imágenes con los mismos factores de reducción
que los formatos con pérdidas. Este tipo de formatos debe utilizarse únicamente
para fotografías o capturas de escenas reales con cámaras digitales.

\index{imagen!vectorial}
Las imágenes vectoriales \textbf{deben} ser empleadas en todo tipo de
diagrama. En ellas no se almacenan píxeles, sino las estructuras geométricas
que componen la figura como círculos (representado por posicion de su centro y
su radio), rectángulos (representados por sus esquinas), líneas, texto, etc. La
mayoría de programas para elaborar este tipo de diagramas, como Inkscape, XFig,
OpenOffice.org Draw, MS Visio, Adobe Illustrator, etc. proveen varios formatos
vectoriales que pueden ser insertados tanto en LaTeX como en OpenOffice.org
Writer (o MS Word). Los formatos más empleados son los llamados metafiles, que
incluyen al WMF, EMF. En LaTeX se utiliza por lo general EPS o PDF, un ejemplo de ello es la Fig. \ref{fig:MMC}.

No debe cometerse el error de generar una imagen vectorial a partir de una
imagen raster, pues una vez realizada la discretización espacial no es posible
reconstruir los elementos geométricos que componen la imagen. Por ello, no
tiene ningún sentido generar un archivo EPS o WMF a partir de una imagen ya
almacenada en BMP, JPG, o PNG, pues lo único que ocurrirá es que se inserta la
figura raster tal cual en la imagen vectorial, sin implicar ninguna ganancia en
la calidad.


\begin{figure}[htb]
  \centering
  \includegraphics[width=0.7\hsize]{MMC}
  \caption{Ejemplo de imagen vectorial.}
  \label{fig:MMC}
\end{figure}

\subsection{Tablas}

También se pueden utilizar tablas, ¡Crearlas es muy fácil!. Puede usar, por ejemplo, el ``creador de tablas online'' \cite{tablesgenerator}.

% Tabla generada con el plugin Excel2Latex
\begin{table}[H] % Importante el H
	\centering
	\caption{Ejemplo de tablas.}
	\begin{tabular}{ccc}
		\hline
		\textbf{Columna 1} & \textbf{Columna 2} & \textbf{Columna 3} \bigstrut\\
		\hline
		$\omega$ & $\nu$ & $\delta$ \bigstrut[t]\\
		$\Phi$ & $\Theta$ & $\varSigma$ \\
		$\Sigma$ & $\epsilon$ & $\psi$ \\
		\hline
	\end{tabular}
	\label{tab:tabla-1}
\end{table}

\subsection{Pseudocódigo}

El pseudocódigo sirve para describir algoritmos de manera detallada y comprensible, sin preocuparse por la sintaxis de un lenguaje de programación específico. Esto permite que los lectores comprendan el flujo lógico del algoritmo y cómo se implementaría en la práctica, sin necesidad de conocer un lenguaje de programación en particular. Además, el pseudocódigo facilita la comunicación de ideas entre escritor y lectores, ya que se centra en los aspectos conceptuales y algorítmicos en lugar de en detalles específicos de implementación.

\begin{algorithm}
\caption{Representación de pseudocódigo de alto nivel de algún algoritmo.}\label{alg:cap}
    \begin{algorithmic}
        \State Calcular un punto inicial factible $u_{0}$.
        \State Definir $\mathcal{V}_{0} = \emptyset $.
        \State $i \gets 0$.
        \While {\textit{true}}
            \If{$\Delta p_{i} = 0$}
                \If{todos los elementos de $\lambda_{i} \geq 0$}
                    \State{$u^{*} \gets u_{i}$}.
                    \State \textbf{break}.
                \Else
                    \State Remover de $\mathcal{V}_{i}$ la restricción con el mínimo $\lambda_{i}$. 
                    \State $u_{i+1} \gets u_{i}$.
                    \State $\mathcal{V}_{i+1} \gets \mathcal{V}_{i}$.
                \EndIf
            \Else
                \State $u_{i+1} \gets u_{i} + \alpha_{i}\Delta p_{i}$.
                \State Crear $\mathcal{V}_{i+1}$ añadiendo la nueva restricción activa.
            \EndIf
            \State $i \gets i+1$.
        \EndWhile
    \end{algorithmic}
\end{algorithm}

\subsection{Referencias bibliográficas}

\index{referencias}\index{BibTeX}
Todo concepto o idea tomado de otros autores contar con la respectiva
referencia. En redacción técnica de ingeniería rara vez se utiliza la cita
textual, así que es necesario reformular las ideas y conceptos con palabras
propias. En ingeniería electrónica se utilizan los formatos de referencia de la
IEEE o la ACM, que son numéricos, encerrados entre paréntesis cuadrados (por
ejemplo, ``En \cite{Davis1963} se propuso un nuevo algoritmo'', o ``En
\cite{ProakisManolakis1998} los autores proponen tomar las ventajas de los
algorimos presentados en \cite{Oppenheim1998,Roberts2005,Haykin2001} por medio
del método de Newton \cite{Burrus1998} conocido en el área de optimización
lineal.''). La referencia es parte de las frases, así que si la frase termina
con la referencia para indicar la idea, ésta debe estar antes del punto final o
demás signos de puntuación: ``La capacidad de memoria también sigue una Ley
similar a la de Moore \cite{Octave}. Los siguientes son los aspectos a tomar en
cuenta en el diseño del sistema \cite{Lindner2002}:''

Se recomienda utilizar BibTeX para administrar las referencias bibliográficas.

